% ladr-macros.tex: notation and environments for the Axler transcription.

% test builds of a single section set their counters only when this is true
\newif\iftestbuild
% Loaded after artbook and macros.tex.  Nothing in artbook.sty is changed.

% ---------------------------------------------------------------- notation
\DeclareRobustCommand{\F}{\mathbb{F}}
\DeclareRobustCommand{\R}{\mathbb{R}}
\DeclareRobustCommand{\C}{\mathbb{C}}
\DeclareRobustCommand{\Q}{\mathbb{Q}}
\DeclareRobustCommand{\Z}{\mathbb{Z}}
\DeclareRobustCommand{\N}{\mathbb{N}}
\DeclareRobustCommand{\Lin}{\mathcal{L}}           % \Lin(V,W)
\DeclareRobustCommand{\Poly}{\mathcal{P}}          % \Poly_m(\F)
\DeclareRobustCommand{\Mat}{\mathcal{M}}           % \Mat(T)
\DeclareMathOperator{\nul}{null}         % \null is a TeX primitive
\DeclareMathOperator{\range}{range}
\DeclareMathOperator{\spn}{span}
\DeclareMathOperator{\adj}{adj}
\DeclareMathOperator{\Real}{Re}
\DeclareMathOperator{\Imag}{Im}
\DeclareMathOperator{\Rank}{rank}
\DeclareMathOperator{\dett}{det}
\DeclareMathOperator{\sgn}{sign}
\renewcommand{\rank}{\operatorname{rank}}
\renewcommand{\tr}{\operatorname{tr}}
\renewcommand{\Re}{\operatorname{Re}}
\renewcommand{\Im}{\operatorname{Im}}
\renewcommand{\dim}{\operatorname{dim}}
\renewcommand{\ip}[2]{\langle #1, #2 \rangle}
\renewcommand{\norm}[1]{\lVert #1 \rVert}
\renewcommand{\abs}[1]{\lvert #1 \rvert}
\newcommand{\ftimes}[2]{\underbrace{#1}_{#2}}  % x\cdots x with "m times" under it

% physics.sty turns \dd, \dv ... into commands; nothing of Axler needs them.
% Axler's l is a plain l: undo the \ell mapping of macros.tex.

% ------------------------------------------------- numbering as in Axler
% Every numbered item (definition, result, example, notation, figure) shares
% one counter per chapter: 1.1, 1.2, ...  Sections are 1A, 1B, ...
\renewcommand\thesection{\thechapter\Alph{section}}
% Axler's subsections are unnumbered (but stay in the contents)
\setcounter{secnumdepth}{1}
\makeatletter
\AtBeginDocument{%
  \@removefromreset{theorem}{section}%
  \@addtoreset{theorem}{chapter}%
  \renewcommand\thetheorem{\thechapter.\arabic{theorem}}%
  % Axler's results carry no name of their own: "Result" in the label column
  \theoremstyle{artbookbox}%
  \newtheorem{result}[theorem]{Result}%
  \newtheorem{notation}[theorem]{Notation}%
  \newtheorem*{result*}{Result}%
  \surroundwithmdframed[style=artbookthm]{result}%
  \surroundwithmdframed[style=artbookthm]{result*}%
  \surroundwithmdframed[style=artbookthm]{notation}%
  % a numbered figure takes its number from the same counter
  \renewcommand\thefigure{\thetheorem}%
}
% \axfigure: a numbered figure in Axler's counter.  Use as
%   \begin{figure}[htbp] ... \axcaption{text}\label{ax:6.9} \end{figure}
\newcommand\axcaption[1]{\refstepcounter{theorem}%
  \addtocounter{figure}{-1}\caption{#1}}

% Fix for artbook's \marginnote: its \raisebox[\ht\z@][\dp\z@]{\box\z@}
% reports a size of zero to \marginpar, so LaTeX never pushes the next note
% below it and consecutive notes overlap.  Set the note as an ordinary
% \marginpar paragraph instead (first baseline on the text line).  The notes
% are placed by marginfix: a note moves down to clear the note above it and
% up when it would run off the foot of the page, so the hand-set offsets
% (\marginnote[offset]{...}) are ignored.
\RequirePackage{marginfix}
\renewcommand\artbook@noteput[2]{\artbook@notestyle #2\par}
% The labels of theorems, proofs, sections and exercise blocks (RESULT 1.3,
% PROOF, SECTION 1A) stand in the same column, but marginfix does not see
% them: reserve their place, so that notes move around them.
% When the page is almost full the box goes to the next page, and a phantom
% left at the foot of this page would keep the notes above it from moving up:
% then there is none.
\newcommand*\ax@labelroom[1]{\par\begingroup
  \skip@=\glueexpr#1\relax \dimen@=\skip@
  \ifdim\dimexpr\pagegoal-\pagetotal\relax<\dimexpr\dimen@+4\baselineskip\relax
    \def\x{\endgroup}%
  \else
    \edef\x{\endgroup\noexpand\marginphantom{\the\dimen@}}%
  \fi\x}
% (added in the preamble, so that it runs before the hook that opens the box)
\@for\ax@env:={definition,definition*,result,result*,notation,example,%
                 example*,remark,remark*,theorem,lemma,proof}\do{%
    \edef\ax@hook{env/\ax@env/before}%
    \expandafter\AddToHook\expandafter{\ax@hook}{\ax@labelroom{\artbook@opt@thmskip+1.4\baselineskip}}}
\makeatother

% ------------------------------------------------- lists inside the text
% (a), (b), ... as Axler sets them; bullets keep the default.
\newlist{parts}{enumerate}{1}
\setlist[parts]{label=(\alph*),leftmargin=*,widest=m,itemsep=1pt,topsep=3pt}
\setlist[itemize]{itemsep=1pt,topsep=3pt,leftmargin=1.2em}
\setlist[enumerate]{itemsep=1pt,topsep=3pt}

% the named property lists of Axler (commutativity, associativity, ...):
%   \begin{axprops} \item[commutativity] text ... \end{axprops}
\newlist{axprops}{description}{1}
\setlist[axprops]{style=nextline,leftmargin=1em,font=\normalfont\bfseries,
  itemsep=2pt,topsep=3pt,parsep=0pt}

% "Exercises 1A": the heading of the exercise block of a section, in the
% manner of a section heading.  A hairline across the page, broken at the
% divider, with a thick bar in the accent colour at its right end; the label
% EXERCISES 1A in the label column and "Exercises" under the bar.
\newif\ifaxafterexercises
\makeatletter
\newcommand*\ax@exrule{%
  \nointerlineskip
  \hbox to\z@{\hskip-\artbook@bodyleft
    \rlap{\artbook@gaprulew{\artbook@fulltw}{\z@}{\artbook@fulltw}{\artbook@bsplit}}%
    \hskip\dimexpr\artbook@bodyleft+\artbook@bodytw-40mm\relax
    {\color{artbook@accent}\rule[-2.6pt]{40mm}{3pt}}\hss}%
  \nobreak\vskip\artbook@opt@secpad\nointerlineskip}
\newcommand\exercisesheading[1]{%
  \par\addvspace{\artbook@opt@secbefore}%
  \Needspace*{6\baselineskip}%
  \global\axafterexercisestrue
  \ax@labelroom{\artbook@opt@secpad+1.4\baselineskip}%
  \begingroup
    \ax@exrule
    \noindent\artbook@leftlabel{\artbook@opt@fontsecnum}{Exercises}{#1}%
    \hfill{\artbook@opt@fontsection Exercises}\par
  \endgroup
  \nobreak\vskip\artbook@opt@secafter\relax
  \@afterindentfalse\@afterheading}

% subsections: a short bar in the accent colour above the title
\renewcommand*\artbook@subsectionstyle{%
  \nointerlineskip
  \hbox{\color{artbook@accent}\rule{10mm}{1.5pt}}%
  \nobreak\vskip 1.8mm\nointerlineskip
  \raggedright\artbook@opt@fontsubsection}
\makeatother

% a short attribution line for the solutions
\newcommand\solnote[2]{\par\smallskip\noindent{\bfseries #1.}\ {\itshape #2}\par}

% a small unnumbered figure in the label column, as Axler's margin figures:
%   \marginfig{<tikz or \includegraphics>}{caption}
% The drawing is set once to measure it, then again with its coordinates
% scaled so that it fills the label column (at most 1.8 times its own
% scale); the text of its labels keeps its size.
\makeatletter
\newsavebox\ax@figbox
\newcommand\marginfig[3][0pt]{%
  \sbox\ax@figbox{#2}%
  \edef\ax@figscale{\fpeval{min(1.8,(\the\marginparwidth)/(\the\wd\ax@figbox))}}%
  \marginnote{\vskip0pt\centering
    {\tikzset{every picture/.append style={scale=\ax@figscale}}#2}\par\smallskip
    {\itshape #3}\par}}
\makeatother

% the quotation that closes some exercise sets
\newcommand\axquote[2]{\par\bigskip\begingroup\raggedleft\small\itshape
  \parbox{.8\linewidth}{\raggedleft #1\par\smallskip\upshape\textemdash #2}\par\endgroup}

% the arrows of the R^2 pictures
\tikzset{vec/.style={-{Stealth[length=2.2mm]},line width=.6pt},
         axis/.style={-{Stealth[length=1.6mm]},line width=.3pt,black!60}}

% solutions end with the black square, as in the manual; \qedhere works in a
% closing display
\AtBeginDocument{\xsimsetup{solution/pre-hook={\pushQED{\qed}},
  solution/post-hook={\popQED}}}

% the small italic remark Axler sets under some exercises
\newcommand\exnote[1]{\par\smallskip{\small\itshape\leftskip=1.5em\relax #1\par}}
\let\axexnote\exnote
\let\exercisenote\exnote

% the text column is narrow: allow a little extra stretch before overfull lines
\setlength\emergencystretch{1.5em}

% same space after the last exercise of every section
\AddToHook{cmd/section/before}{\par\addvspace{\bigskipamount}}
% the section after an exercise block starts on a new page (the quotation
% that closes some exercise blocks stays with them)
\AddToHook{cmd/section/before}{%
  \ifaxafterexercises\clearpage\global\axafterexercisesfalse\fi}
\AddToHook{cmd/chapter/before}{\global\axafterexercisesfalse}
\makeatletter
\AddToHook{cmd/section/before}{\ax@labelroom{3.8\baselineskip+\artbook@opt@secpad}}
\makeatother
\AddToHook{cmd/part/before}{\global\axafterexercisesfalse}
\allowdisplaybreaks[1]

% artbook's contents lists only section-level rows (chapters are headers),
% so each group of solutions gets its own section-level entry
\renewcommand{\printsolutionsbysection}{%
  \def\lastchapter{}\def\lastsection{}%
  \ForEachUsedExerciseByType{%
    \edef\currentchapter{\ExercisePropertyGet{##1}{##2}{chapter-value}}%
    \edef\currentsection{\ExercisePropertyGet{##1}{##2}{section-value}}%
    \ifx\lastchapter\currentchapter\else
      \section*{Chapter~\ExercisePropertyGet{##1}{##2}{chapter}}%
    \fi
    \ifx\lastsection\currentsection\else
      \subsection*{\S\ExercisePropertyGet{##1}{##2}{section}}%
      \addcontentsline{toc}{section}{Solutions to
        \ExercisePropertyGet{##1}{##2}{section}}%
    \fi
    \XSIMprint{solution}{##1}{##2}%
    \let\lastchapter\currentchapter
    \let\lastsection\currentsection}}

% artbook counts only numbered chapter heads but every chapter line in the
% .toc, so unnumbered front-matter chapters shifted every chapter's mini
% contents by five.  Count the starred heads too.
\makeatletter
\def\@makeschapterhead#1{\global\advance\artbook@chapcnt\@ne
  \chaptermark{#1}\artbook@chapterpage{#1}{0}}
\makeatother
